\documentclass[11pt,a4paper]{article}
\usepackage[T1]{fontenc}
\usepackage[utf8]{inputenc}
\usepackage{lmodern}
\usepackage[margin=0.85in]{geometry}
\usepackage{parskip}
\setlength{\parskip}{0.55em}
\usepackage{hyperref}
\hypersetup{hidelinks}
\pagestyle{empty}

\begin{document}

\noindent
Almaz Sapargali (Corresponding Author)\\
LLP Digit Alem, Almaty, Kazakhstan\\
Al-Farabi Kazakh National University, Almaty, Kazakhstan\\
\href{mailto:almaz.sapargali2001@gmail.com}{almaz.sapargali2001@gmail.com}

\vspace{1em}
\noindent\today

\vspace{1em}
\noindent
To the Editorial Office / Editor-in-Chief\\
\textit{AI} (MDPI), ISSN 2673-2688

\vspace{1.2em}
\noindent Dear Editor,

\vspace{0.6em}
We are pleased to submit our manuscript, ``\textbf{The Fidelity--Utility Paradox
in Surrogate-Based Reinforcement Learning: A Role-Separated Hybrid Methodology in
an HVAC Control Case Study},'' for consideration as an original research Article in
\textit{AI}.

Deep reinforcement learning (DRL) for building HVAC control almost universally
relies on fast learned surrogates as training environments, under the tacit
assumption that a more predictively accurate surrogate is a better environment in
which to train a controller. To our knowledge, our study is the first to test
this assumption directly and, under a pre-specified, version-locked protocol, to report that it
can fail. On the BOPTEST \texttt{bestest\_air} benchmark, a calibrated grey-box
physical twin that is more than twice as accurate as a coarse black-box surrogate
(24-hour rollout RMSE $0.644$ vs.\ $1.557\,^{\circ}$C) nevertheless yields
controllers that collapse on every seed on the live simulator, whereas the weaker
surrogate trains a usable controller. We show that the operative variable is the
surrogate's temporal resolution (its coarse per-step increment) rather than
model class or predictive accuracy, and we resolve the paradox with a
role-separating hybrid: black-box rollouts combined with a frozen physics twin
used only as a per-step model-disagreement reward censor, which recovers robust
control at roughly $85\times$ live-simulator throughput. A pre-specified transferability
study then delineates a precise, component-level boundary: the inverse-calibration
pipeline transfers across a hydronic family of buildings, whereas a frozen policy
does not.

We believe this work is a strong fit for \textit{AI} because it speaks directly to
how learned world models and surrogate environments should be evaluated and used
for control. HVAC serves as the concrete case study for a broader
surrogate-based-RL hypothesis: rather than proposing yet another controller, the
paper provides a corrected, quantified account of \emph{why} surrogate-trained
controllers succeed or fail, together with a reusable role-separating recipe, of
practical significance for the growing literature on physics-informed digital twins
and safe RL for buildings.

This manuscript is original, has not been published previously, and is not under
consideration elsewhere. All authors have read and approved the submitted version
and agree to its submission to \textit{AI}. The authors declare no conflicts of
interest, and this simulation-only study involves no human or animal subjects. This
research was funded by the Committee of Science of the Ministry of Science and
Higher Education of the Republic of Kazakhstan (Grant No.~AP23488794). Thank you for
considering our submission; we look forward to the reviewers' comments.

\vspace{1.2em}
\noindent Sincerely,

\vspace{0.4em}
\noindent
Almaz Sapargali\\
on behalf of all co-authors\\
(Aksultan Mukhanbet, Serik Aibagarov, Beimbet Daribayev, Shona Shinassylov, Paulo Trigo)

\end{document}
